\chapter{Cicero or Christ?}
"The name of Jesus can still remove distractions from the minds of
men."--\textit{Origen}.


Shall men believe in a future life because of Christ's return from
the grave?  Is his established resurrection at Jerusalem the
climacteric proof for immortality?  The problem is inescapable.
Every man is himself a judge; before every man the accumulated
evidence passes; for every man it is doomsday when he stands at the
point of decision.

In his sore perplexity Quintus says to himself that night, when he
has returned from his interview with the disciple John: "My soul is
like a traveler who halts at the point where two roads meet.  Great
issues depend upon his choice.  But while he hesitates may the
immortals, who watch over the destinies of men, guide his feet
aright."

Clearly defined are the alternatives before the Roman soldier.  On
the one hand are his ancestral beliefs, long established and deeply
cherished by the nation.  Nor does any man quickly toss aside the
faith of his fathers.  If belief is waning in the primitive
mythologies, and if the social life of the Empire is moved by
unrest and despair, the problem is to find a greater satisfaction.
There have been spoken many beautiful words by the Roman scholars
which are sweet premonitions of immortality.  Does not Quintus
remember that Cicero likens to heaven a port prepared, and prays
that he may sail thither with full-spread sails?  And if the gifted
Cicero has just gone tragically out of life, let it be hoped that
he has reached the harbor.

But on the other hand are the challenging and captivating words of
Christ.  Had he only spoken of the future life as an enthusiastic
Teacher, and then had passed to the perpetual slumber of the grave
like other philosophers of the time, he would be remembered long.
But, when he had spoken his words concerning immortality, he had
added, "I myself shall surely come back again."  From the evidence
which Quintus has heard in Jerusalem he has now fulfilled his
prediction.  He has put to scorn the fidelity of the Roman
sentinels at the tomb of Joseph; he has reversed the laws of
nature; he has appeared again, in unique proof that there is to be
a resurrection of the dead.  Wide is the difference between Cicero
and the Christ.  The one has spoken a mere opinion, so beautiful in
its phrase that it shall pass down into the future literature of
men.  The other has spoken a revelation, and then has returned to
prove that revelation true.  Which shall it be--Cicero or the
Christ?

But to accept the Jewish Teacher means earthly loss.  As he keeps
guard with himself through the night hours Quintus is wondering if
he shall incur the hostility of his father Marcus and shall be
forced to sacrifice his estates on the Palatine.  He fancies also
the grief of the fair Lucretia when she learns that he has chosen
an alien faith.  And he remembers, further, that in the choice of
the Christus he is joining a company on whom the Eastern world is
already casting its withering contempt.  Cicero or the Christus.
Which shall it be?


There are no struggles like the night wrestlings of the soul in
matters of religion.  What words can measure the divers arguments,
the opposing considerations, the conflicting emotions that shape
human choice?  Quintus stands at the point where soon--in the
progress of the new faith--Saul from Tarsus, Clement of Rome, and
so many more of the great spirits of that first era are to stand.
The wrestlings of the night!  Then foul demons are abroad; and then
God's good angels are descending the ladders of the sky.

Soon comes a great moment.  While the soul of Quintus is in wild
commotion, there falls upon him a mighty force which is not of
earth.  Coming he knows not whence, but not invading the department
of his will, it impels him to the Christ.  Transformed is this
Roman knight, who has been taught the doctrines of the Latin cult,
and whose nation can only feel disdain for a Galilaean who proposes
to revolutionize the ages.  The words of the augur at Brundisium
are having in truth a strange fulfillment.

As if the Man were present on whom he had looked in the Porch of
Solomon.  Quintus speaks his choice for the long eternities:

"Happen what may, I take thee, O Christus, for my Lord and Master.
I sacrifice my Roman knighthood for thee, if it shall be required.
I choose thee, because thou hast risen from the dead and hast
proven that there is another life for men."

Not Cicero, but Christ!  The Roman knight has made the great
decision.